\chapter{主动性与持续追求}

\section{期望与主动性}

Nous Wave 可以形成期望条件。期望条件描述 Subject 希望成立、停止、持续或避免的状态。

Heptalogos 使用主动追求（Pursuit）表示已经被行为侧采纳的持续主动追求。Pursuit 引用期望条件的精确修订，并保留采纳原因、当前状态和来源。

\section{持续追求}

主动追求可以跨越多次 Reaction、会话和时间窗口。它能够产生新的 CognitiveOpportunity，影响注意和行为判断，并参与资源竞争。

期望条件本身不触发外部效果。Heptalogos 根据情境、政策、现有承诺、能力、时间和资源决定是否形成或继续 Pursuit。

期望条件修订后，引用旧修订的主动追求需要重新核验。已经发生的行为提交、承诺和外部效果保持历史真实性。

\section{相关状态}

行为提议（BehaviorIntent）描述一次行为判断中的候选行为。行动计划（ActionPlan）或交互计划（InteractionPlan）描述实现方式。承诺（Commitment）与义务（Obligation）描述已经承担的约束。主动追求表示跨多次行为判断延续的追求状态。

这些对象分别维护身份和生命周期。
